\chapter{Desafíos en las aplicaciones distribuidas}
Las aplicaciones distribuidas son fundamentales en la arquitectura de sistemas modernos, ofreciendo ventajas significativas en términos de escalabilidad, flexibilidad y disponibilidad. Sin embargo, estas ventajas vienen acompañadas de desafíos complejos derivados de la naturaleza intrínseca de la distribución de componentes a lo largo de redes, posiblemente extendidas a través de vastas geografías. Uno de los retos más significativos es asegurar la consistencia de los datos y el estado del sistema a través de múltiples nodos, lo cual es crítico para mantener la integridad y la precisión operacional. Además, las aplicaciones distribuidas deben gestionar la latencia de la red y las potenciales particiones de red, que pueden afectar negativamente la experiencia del usuario y la eficacia del sistema.

Otro aspecto desafiante es la tolerancia a fallos y la recuperación ante errores en entornos distribuidos. A medida que el número de componentes en un sistema distribuido aumenta, también lo hace la probabilidad de fallos en algún punto del sistema. Implementar estrategias efectivas de detección de fallos, recuperación y redundancia es esencial para asegurar la alta disponibilidad y la continuidad del servicio. Estos desafíos requieren un diseño meticuloso y consideraciones operativas que abarquen desde la elección de la tecnología hasta la implementación de políticas de seguridad robustas para proteger los datos transmitidos a través de redes potencialmente inseguras. Así, el desarrollo y mantenimiento de aplicaciones distribuidas exige un enfoque holístico y detallado que equilibre eficazmente la funcionalidad, el rendimiento y la seguridad.
